\setchapterimage{}
\setchapterintro{}
\chapter{概率空间与随机事件}

\section{随机试验和事件}
\label{sec:random-experiment-and-event}

\subsection*{随机试验、样本空间与事件}

在固定条件下必定出现唯一结果的现象称为确定性现象；在固定条件下可能出现两个或两个以上结果的现象称为随机现象。概率论与数理统计研究随机现象的数量规律。

\begin{mydefinition}{随机试验}{random-experiment}
在概率论中，将满足下列条件的试验称为随机试验，简称试验：
\begin{enumerate}
\item 试验进行前，无法确切知道一定会出现哪个结果；
\item 已知试验所有可能的结果；
\item 试验结束后，能够明确知道哪个试验结果出现。
\end{enumerate}
\end{mydefinition}

\begin{mydefinition}{样本点与样本空间}{sample-space}
\begin{enumerate}
\item 一个随机试验的任意一个可能结果称为一个\textbf{样本点}（sample），记为 $\omega$；
\item 该随机试验所有可能结果组成的集合称为该试验的\textbf{样本空间}（sample space），记为 $\Omega$。
\end{enumerate}
\end{mydefinition}

样本点和样本空间的关系为
\[
\omega\in\Omega.
\]
样本空间 $\Omega$ 可以是有限集、无限集，也可以是以函数为元素的集合。用集合语言描述试验结果具有抽象性和普遍性，便于刻画不同随机试验的共同结构。

\begin{mydefinition}{事件}{event}
\begin{enumerate}
\item 样本空间 $\Omega$ 的一个子集称为一个\textbf{事件}，通常用大写字母 $A,B,C,\ldots$ 表示；
\item 对事件 $A$，若试验结果 $\omega\in A$，则称事件 $A$ 发生；否则称事件 $A$ 不发生。
\end{enumerate}
\end{mydefinition}

\textbf{注:(基本事件)}样本空间 $\Omega$ 包含所有试验结果，故称为\textbf{必然事件}；空集 $\varnothing$ 不含任何样本点，故称为\textbf{不可能事件}。

下面用集合的关系与运算刻画事件的关系与运算。

\begin{enumerate}
\item \textbf{事件的包含与相等}

若 $A\subseteq B$，则称事件 $B$ 包含事件 $A$，或称事件 $A$ 包含在事件 $B$ 中。其概率意义是：若事件 $A$ 发生，则事件 $B$ 必发生。

若 $A\subseteq B$ 且 $B\subseteq A$，即 $A=B$，则称事件 $A$ 与事件 $B$ 相等；其概率意义是二者同时发生、同时不发生。

\item \textbf{事件的并}

集合 $A\cup B$ 称为事件 $A$ 与 $B$ 的并事件，其概率意义是 $A$ 与 $B$ 至少有一个发生。更一般地，
\[
\bigcup_{i=1}^{n} A_i
=\left\{\omega\in\Omega\mathrel{\big|}\exists\, i\in\{1,\ldots,n\},\ \omega\in A_i\right\},
\]
表示 $n$ 个事件中至少有一个发生；而
\[
\bigcup_{i=1}^{\infty} A_i
=\left\{\omega\in\Omega\mathrel{\big|}\exists\, i\geq 1,\ \omega\in A_i\right\},
\]
表示可数个事件中至少有一个发生。

\item \textbf{事件的交（积）}

集合 $A\cap B$ 称为事件 $A$ 与 $B$ 的交（积）事件，其概率意义是 $A$ 与 $B$ 都发生。类似地，
\[
\bigcap_{i=1}^{n} A_i
=\left\{\omega\in\Omega\mathrel{\big|}\forall\, i\in\{1,\ldots,n\},\ \omega\in A_i\right\},
\]
表示 $n$ 个事件都发生；
\[
\bigcap_{i=1}^{\infty} A_i
=\left\{\omega\in\Omega\mathrel{\big|}\forall\, i\geq 1,\ \omega\in A_i\right\},
\]
表示可数个事件都发生。为书写方便，常记 $AB:=A\cap B$。

\item \textbf{事件的差}

集合
\[
A\setminus B=\left\{\omega\in\Omega\mathrel{\big|}\omega\in A,\ \omega\notin B\right\}
\]
称为事件 $A$ 与 $B$ 的差，其概率意义是 $A$ 发生而 $B$ 不发生。

\item \textbf{互斥（互不相容）事件}

若 $A\cap B=\varnothing$，则称事件 $A$ 与 $B$ 互斥或互不相容，其概率意义是二者不能同时发生。若对一列事件 $A_n\ (n\geq1)$，任意 $i\ne j$ 都有 $A_i\cap A_j=\varnothing$，则称该列事件两两互斥。

\item \textbf{补事件}

事件 $A$ 的补集
\[
A^c=\left\{\omega\in\Omega\mathrel{\big|}\omega\notin A\right\}
\]
称为 $A$ 的补事件，其概率意义是 $A$ 不发生。并且
\[
A\setminus B=A\cap B^c.
\]
事件运算约定先进行补运算，再进行交运算，最后进行并运算或差运算。

\item \textbf{完备事件组}

若事件 $A_1,A_2,\ldots$ 两两互斥，且
\[
\bigcup_{i=1}^{\infty}A_i=\Omega,
\]
则称 $A_1,A_2,\ldots$ 构成一个\textbf{完备事件组}，或称为 $\Omega$ 的一个可数划分。
\end{enumerate}

事件的运算性质就是集合的运算性质。对任意事件 $A,B,C$，有
\[
A\cup B=B\cup A,\qquad A\cap B=B\cap A,
\]
\[
(A\cup B)\cup C=A\cup(B\cup C),\qquad
(A\cap B)\cap C=A\cap(B\cap C),
\]
\[
A\cup(B\cap C)=(A\cup B)\cap(A\cup C),\qquad
A\cap(B\cup C)=(A\cap B)\cup(A\cap C).
\]
此外，德摩根律为
\[
\left(\bigcup_{i=1}^{n}A_i\right)^c=\bigcap_{i=1}^{n}A_i^c,
\qquad
\left(\bigcap_{i=1}^{n}A_i\right)^c=\bigcup_{i=1}^{n}A_i^c,
\]
以及
\[
\left(\bigcup_{i=1}^{\infty}A_i\right)^c=\bigcap_{i=1}^{\infty}A_i^c,
\qquad
\left(\bigcap_{i=1}^{\infty}A_i\right)^c=\bigcup_{i=1}^{\infty}A_i^c.
\]

\subsection*{古典概率模型（高中内容）}

对事件 $A\subseteq\Omega$，用 $P(A)$ 表示事件 $A$ 发生的概率。

\begin{mydefinition}{古典概率模型}{classical-probability-model}
若某随机试验的样本空间 $\Omega$ 只含有限个样本点，且每个样本点出现的概率相等，则称该随机试验为\textbf{古典概率模型}，简称\textbf{古典概型}。对任意事件 $A\subseteq\Omega$，记 $A$ 中元素个数为 $n(A)$，则
\[
P(A)=\dfrac{n(A)}{n(\Omega)}.
\]
\end{mydefinition}

由古典概型的定义可知，计算 $P(A)$ 的关键在于计算 $n(A)$ 和 $n(\Omega)$，因此需要使用计数方法。

\begin{proposition}{计数的乘法原理}{multiplication-principle}
若一个过程由 $r$ 个试验构成，第 $i$ 个试验有 $n_i$ 种可能结果，则该过程共有
\[
n_1n_2\cdots n_r
\]
种可能结果。
\end{proposition}
\begin{example}
	\begin{enumerate}
		\item 有放回有序抽取：从 $n$ 个不同的球中有放回地每次抽取一个，依次抽取 $m$ 次并按顺序排成一列，可得到 $n^m$ 个不同排列。
		\item 无放回有序抽取：从 $n$ 个不同的球中无放回地抽取 $m$ 个并按顺序排成一列，可得到
		\[
		A_n^m=n(n-1)\cdots(n-m+1)=\dfrac{n!}{(n-m)!}
		\]
		个不同排列。
		\item 圆排列数：$n$个人坐一圈，有$(n-1)!$种排列方式.
		\item 组合数：从 $n$ 个不同的球中无放回地抽取 $m$ 个并放在一起，不计顺序，可得到
		\[
		C_n^m=\binom{n}{m}=\dfrac{n!}{(n-m)!m!}
		\]
		个不同组合。
	\end{enumerate}
\end{example}

\begin{proposition}{多组分配}{multinomial-count}
将 $n$ 个不同元素分成 $r$ 个组，其中第 $i$ 组有 $n_i$ 个元素，组内不计顺序，且
\[
n_1+n_2+\cdots+n_r=n,
\]
则不同分法共有
\[
\dfrac{n!}{n_1!n_2!\cdots n_r!}
\]
种。
\end{proposition}

\begin{proof}
依次选出各组元素，不同分法数为
\[
\binom{n}{n_1}
\binom{n-n_1}{n_2}
\binom{n-n_1-n_2}{n_3}
\cdots
\binom{n-n_1-\cdots-n_{r-1}}{n_r}.
\]
将各组合数相乘并约去中间阶乘，得到
\[
\dfrac{n!}{n_1!n_2!\cdots n_r!}.
\]
\end{proof}

\begin{mycorollary}{多项式定理}{multinomial-theorem}
\[
(x_1+x_2+\cdots+x_r)^n
=\sum_{n_1+n_2+\cdots+n_r=n}
\dfrac{n!}{n_1!n_2!\cdots n_r!}
x_1^{n_1}x_2^{n_2}\cdots x_r^{n_r}.
\]
其中
\[
\binom{n}{n_1,n_2,\ldots,n_r}
:=\dfrac{n!}{n_1!n_2!\cdots n_r!}
\]
称为多项式系数。
\end{mycorollary}
\newpage
\begin{proposition}{正整数解的个数}{positive-integer-solutions}
设 $n\geq1$，方程
\[
x_1+x_2+\cdots+x_r=n,\qquad x_i>0\quad(i=1,2,\ldots,r)
\]
有
\[
\binom{n-1}{r-1}
\]
个不同的正整数解。
\end{proposition}

\begin{proof}
将 $n$ 个 $0$ 排成一行，从其中的 $n-1$ 个间隔中选出 $r-1$ 个插入隔板。令 $x_1$ 为第一个隔板左侧 $0$ 的个数，$x_2$ 为前两个隔板之间 $0$ 的个数，依此类推，$x_r$ 为最后一个隔板右侧 $0$ 的个数。

每一种隔板插法恰对应方程的一个正整数解，故不同解的个数为
\[
\binom{n-1}{r-1}.
\]
\end{proof}

\begin{proposition}{非负整数解的个数}{nonnegative-integer-solutions}
设 $n\geq1$，方程
\[
x_1+x_2+\cdots+x_r=n,\qquad x_i\geq0\quad(i=1,2,\ldots,r)
\]
有
\[
\binom{n+r-1}{r-1}
\]
个不同的非负整数解。
\end{proposition}

\begin{proof}
令 $y_i=x_i+1\ (i=1,2,\ldots,r)$。则 $x_1,x_2,\ldots,x_r$ 为满足原方程的非负整数解，当且仅当 $y_1,y_2,\ldots,y_r$ 为满足
\[
y_1+y_2+\cdots+y_r=n+r,\qquad y_i>0
\]
的正整数解。由正整数解的个数公式，所求解数为
\[
\binom{n+r-1}{r-1}.
\]
\end{proof}

\subsection*{几何概率模型（高中内容）}

\begin{mydefinition}{几何概率模型}{geometric-probability-model}
设随机试验的样本空间 $\Omega$ 是 $\mathbb{R}^d$ 中一个具有体积 $m(\Omega)$ 的子集，且
\[
0<m(\Omega)<+\infty.
\]
若事件 $A\subseteq\Omega$ 具有体积 $m(A)$，并规定
\[
P(A)=\dfrac{m(A)}{m(\Omega)},
\]
则称该随机试验为\textbf{几何概率模型}，简称\textbf{几何概型}。
\end{mydefinition}

\textbf{注：}几何概型也是一种等可能模型：对具有相同体积（在平面上即相同面积）的事件，它们发生的可能性相同。

\vspace{2\baselineskip}

\begin{example}[贝特朗奇论]
随机地从半径为 $r$ 的圆中取一根弦，求该弦长度大于圆内接正三角形边长 $\sqrt{3}r$ 的概率。设事件 $A$ 表示“弦长大于 $\sqrt{3}r$”。

“随机取弦”的说法本身并未说明等可能的参照标准，因此可以得到不同的概率。
\begin{enumerate}
\item 若以弦的中点 $M$ 确定一根弦，并规定 $M$ 在大圆内均匀取值，则弦长大于 $\sqrt{3}r$ 当且仅当 $M$ 落在半径为 $r/2$ 的同心小圆内。因此
\[
P(A)=\dfrac{\pi(r/2)^2}{\pi r^2}=\dfrac14.
\]

\item 若固定弦的一个端点，并规定弦与该端点处切线的夹角 $\theta$ 在 $[0,\pi]$ 上均匀取值，则
\[
A=\left(\dfrac{\pi}{3},\dfrac{2\pi}{3}\right),
\qquad
P(A)=\dfrac{\pi/3}{\pi}=\dfrac13.
\]

\item 若以弦到圆心的距离 $D$ 确定一根弦，并规定 $D$ 在 $[0,r]$ 上均匀取值，则 $A=\{0<D<r/2\}$，从而
\[
P(A)=\dfrac{r/2}{r}=\dfrac12.
\]
\end{enumerate}

三种理解得到不同答案的原因在于：它们分别选择了不同的样本空间和等可能标准。求概率前，必须先明确随机试验的样本空间以及概率模型。
\end{example}

\begin{example}[布丰投针实验]
在平面上画出间距为 $d$ 的一组平行直线。随机投掷一根长度为 $l$ 的针，设
\[
0<l\leq d.
\]
记针与任一直线相交为事件 $A$。令 $X$ 表示针的中点到最近平行线的距离，$\Theta$ 表示针与平行线的夹角。按照几何概型，取
\[
0\leq X\leq \dfrac d2,
\qquad
0\leq \Theta\leq \dfrac{\pi}{2},
\]
并把 $(X,\Theta)$ 看作在该矩形内等可能取值。

针与直线相交当且仅当
\[
X\leq \dfrac l2\sin\Theta.
\]
因此，事件 $A$ 对应的面积为
\[
\int_0^{\pi/2}\dfrac l2\sin\theta\,\mathrm{d}\theta
=\dfrac l2,
\]
而样本空间的面积为 $\dfrac d2\cdot\dfrac{\pi}{2}=\dfrac{d\pi}{4}$。故
\[
P(A)=\dfrac{l/2}{d\pi/4}=\dfrac{2l}{\pi d}.
\]

若独立投针 $N$ 次，其中有 $H$ 次与直线相交，则当 $N$ 较大时，频率 $H/N$ 可近似看作 $P(A)$，从而得到 $\pi$ 的实验估计式
\[
\pi\approx\dfrac{2lN}{dH}.
\]
\end{example}

\vspace{2\baselineskip}

\section{可测空间和概率空间}

古典概型和几何概型都是等可能模型，但并非所有随机试验都能归结为这两类模型。贝特朗奇论还表明，必须用更严格的形式明确样本空间和概率的含义，才能避免由语义歧义带来的不同答案。

\subsection*{面积和频率}

设 $A\subseteq\mathbb{R}^d$，用 $\mu(A)$ 表示 $A$ 的面积、长度或体积。直观上，$\mu(\cdot)$ 具有下列性质：
\begin{enumerate}
\item $\mu(A)\geq0$；
\item $\mu(\varnothing)=0$；
\item 若 $A\subseteq B$，且 $0\leq\mu(A),\mu(B)<+\infty$，则
\[
\mu(B\setminus A)=\mu(B)-\mu(A);
\]
\item 若 $A_1,A_2,\ldots$ 两两互斥，则
\[
\mu\left(\bigcup_{n=1}^{\infty}A_n\right)
=\sum_{n=1}^{\infty}\mu(A_n).
\]
\end{enumerate}

日常生活中也常以频率描述事件发生的可能性。频率的性质将启发我们给出概率的公理化定义。

\begin{mydefinition}{频率}{frequency}
在 $N$ 次重复试验中，若事件 $A$ 出现 $N(A)$ 次，则比值
\[
f_N(A)=\dfrac{N(A)}{N}
\]
称为事件 $A$ 发生的\textbf{频率}。
\end{mydefinition}

\begin{proposition}{频率的性质}{frequency-properties}
对任意事件 $A$，有
\[
0\leq f_N(A)\leq1,
\qquad
f_N(\Omega)=1,
\qquad
f_N(\varnothing)=0.
\]
并且
\[
f_N(A^c)=1-f_N(A).
\]
若 $A_1,A_2,\ldots,A_r$ 两两互斥，则
\[
f_N\left(\bigcup_{i=1}^{r}A_i\right)
=\sum_{i=1}^{r}f_N(A_i).
\]
\end{proposition}

\begin{proof}
由 $0\leq N(A)\leq N$，两端除以 $N$ 即得 $0\leq f_N(A)\leq1$。每次试验中 $\Omega$ 都发生而 $\varnothing$ 都不发生，故 $N(\Omega)=N$、$N(\varnothing)=0$。

若 $A_1,\ldots,A_r$ 两两互斥，则每次试验至多落入其中一个事件，因此
\[
N\left(\bigcup_{i=1}^{r}A_i\right)=\sum_{i=1}^{r}N(A_i).
\]
两端除以 $N$ 即得结论。

又因为在每次试验中，$A$ 与 $A^c$ 恰有一个发生，故 $N(A^c)=N-N(A)$，从而
\[
f_N(A^c)=\dfrac{N-N(A)}{N}=1-f_N(A).
\]
\end{proof}

频率反映了事件在重复试验中发生的频繁程度，并直观地反映事件在一次试验中发生的可能性大小。

\subsection*{可测空间}

由面积和频率的性质可见，面积 $\mu(\cdot)$、频率 $f_N(\cdot)$ 以及概率 $P(\cdot)$ 都是以集合为自变量的函数。为了说明哪些子集可以作为事件，需要先规定一个对补集和可数并封闭的集合类。

\begin{mydefinition}{$\sigma$-域与可测空间}{sigma-algebra}
设 $\Omega$ 是非空集合，$\mathscr{F}$ 是由 $\Omega$ 的一些子集构成的集合类。若 $\mathscr{F}$ 满足：
\begin{enumerate}
\item $\Omega\in\mathscr{F}$；
\item 若 $A\in\mathscr{F}$，则 $A^c\in\mathscr{F}$；
\item 若 $A_n\in\mathscr{F}$（$n\geq1$），则
\[
\bigcup_{n=1}^{\infty}A_n\in\mathscr{F},
\]
\end{enumerate}
则称 $\mathscr{F}$ 为 $\Omega$ 的一个\textbf{$\sigma$-域}或\textbf{事件域}，在分析学中我们称之为\textbf{$\sigma$-代数}，$\mathscr{F}$ 中的元素称为事件或可测集，二元组 $(\Omega,\mathscr{F})$ 称为\textbf{可测空间}。
\end{mydefinition}

\begin{proposition}{$\sigma$-域的基本封闭性}{sigma-algebra-closure}
若 $\mathscr{F}$ 是 $\Omega$ 的一个 $\sigma$-域，则：
\begin{enumerate}
\item $\varnothing\in\mathscr{F}$；
\item 若 $A_n\in\mathscr{F}$（$n\geq1$），则
\[
\bigcap_{n=1}^{\infty}A_n\in\mathscr{F};
\]
\item 若 $A_1,\ldots,A_k\in\mathscr{F}$，则有限并与有限交都属于 $\mathscr{F}$；
\item 若 $A,B\in\mathscr{F}$，则 $A\setminus B\in\mathscr{F}$。
\end{enumerate}
\end{proposition}

\begin{proof}
由 $\Omega\in\mathscr{F}$ 及补集封闭性，得 $\varnothing=\Omega^c\in\mathscr{F}$。若 $A_n\in\mathscr{F}$，则 $A_n^c\in\mathscr{F}$，从而
\[
\bigcap_{n=1}^{\infty}A_n
=\left(\bigcup_{n=1}^{\infty}A_n^c\right)^c
\in\mathscr{F}.
\]
有限并可视为可数并中从第 $k+1$ 项起均取空集的情形；有限交由德摩根律得到。最后，$B^c\in\mathscr{F}$，故
\[
A\setminus B=A\cap B^c\in\mathscr{F}.
\]
\end{proof}

\begin{example}[一些 $\sigma$-代数的例子]
\begin{enumerate}
\item $\{\varnothing,\Omega\}$ 是 $\Omega$ 上所含事件最少的 $\sigma$-域；
\item 对任意 $A\subseteq\Omega$，集合类
\[
\{\varnothing,\Omega,A,A^c\}
\]
是包含 $A$ 的一个 最小的$\sigma$-域；
\item $\mathcal{P}(\Omega):=\{A\mid A\subseteq\Omega\}$ 是 $\Omega$ 的所有子集构成的 $\sigma$-域。
\end{enumerate}
\end{example}

\textbf{注：}样本点 $\omega\in\Omega$ 并不必然意味着单点集 $\{\omega\}$ 属于 $\mathscr{F}$。因此，只有在固定可测空间 $(\Omega,\mathscr{F})$ 后，才能判断某个子集是否为事件。

\begin{proposition}{$\sigma$-域的交}{intersection-of-sigma-algebras}
设 $I$ 是一个指标集。若对任意 $\alpha\in I$，$\mathscr{F}_\alpha$ 都是 $\Omega$ 上的 $\sigma$-域，则
\[
\mathscr{F}=\bigcap_{\alpha\in I}\mathscr{F}_\alpha
\]
也是 $\Omega$ 上的 $\sigma$-域。
\end{proposition}

\begin{proof}
对每个 $\alpha\in I$，都有 $\Omega\in\mathscr{F}_\alpha$，故 $\Omega\in\mathscr{F}$。若 $A\in\mathscr{F}$，则 $A\in\mathscr{F}_\alpha$ 对所有 $\alpha\in I$ 成立，因而 $A^c\in\mathscr{F}_\alpha$ 对所有 $\alpha$ 成立，即 $A^c\in\mathscr{F}$。同理，若 $A_n\in\mathscr{F}$（$n\geq1$），则 $\bigcup_{n=1}^{\infty}A_n\in\mathscr{F}_\alpha$ 对所有 $\alpha$ 成立，故该可数并属于 $\mathscr{F}$。
\end{proof}

\begin{mytheorem}{由集合类生成的 $\sigma$-域}{generated-sigma-algebra}
设 $\Omega$ 是非空集合，$\mathscr{E}$ 是由 $\Omega$ 的一些子集构成的非空集合类。则存在唯一的 $\sigma$-域 $\mathscr{F}$ 满足：
\begin{enumerate}
\item $\mathscr{E}\subseteq\mathscr{F}$；
\item 对任意包含 $\mathscr{E}$ 的 $\sigma$-域 $\mathcal{M}$，都有 $\mathscr{F}\subseteq\mathcal{M}$。
\end{enumerate}
此时称 $\mathscr{F}$ 为由 $\mathscr{E}$ \textbf{生成的 $\sigma$-域}，记为 $\sigma(\mathscr{E})$。
\end{mytheorem}

\begin{proof}
令
\[
\mathscr{F}
=\bigcap_{\substack{\mathscr{M}\text{ 是 }\Omega\text{ 上的 }\sigma\text{-域}\\ \mathscr{E}\subseteq\mathscr{M}}}\mathscr{M}.
\]
因为 $\mathcal{P}(\Omega)$ 是包含 $\mathscr{E}$ 的 $\sigma$-域，故上式中的交非空。由上一命题，$\mathscr{F}$ 是 $\sigma$-域；由定义，$\mathscr{E}\subseteq\mathscr{F}$，并且 $\mathscr{F}$ 包含于任意包含 $\mathscr{E}$ 的 $\sigma$-域。唯一性由这一最小性立刻得到。
\end{proof}

\begin{mydefinition}{Borel $\sigma$-域}{borel-sigma-algebra}
记
\[
\mathscr{E}_{\mathbb{R}}=\{(a,b]\mid a,b\in\mathbb{R}\}.
\]
称 $\sigma(\mathscr{E}_{\mathbb{R}})$ 为 $\mathbb{R}$ 上的\textbf{Borel $\sigma$-域}，记为 $\mathscr{B}(\mathbb{R})$。

类似地，记
\[
\mathscr{E}_{\mathbb{R}^d}
=\left\{(a_1,b_1]\times\cdots\times(a_d,b_d]\mid a_i,b_i\in\mathbb{R},\ 1\leq i\leq d\right\}.
\]
称 $\sigma(\mathscr{E}_{\mathbb{R}^d})$ 为 $\mathbb{R}^d$ 上的 Borel $\sigma$-域，记为 $\mathscr{B}(\mathbb{R}^d)$。
\end{mydefinition}

\begin{proposition}{常用集合属于 $\mathscr{B}(\mathbb{R})$}{common-borel-sets}
对任意 $a<b$，开区间 $(a,b)$、闭区间 $[a,b]$、半开半闭区间 $[a,b)$ 以及单点集 $\{b\}$ 都属于 $\mathscr{B}(\mathbb{R})$。
\end{proposition}

\begin{proof}
对任意 $k\geq1$，$(a,b-1/k]\in\mathscr{B}(\mathbb{R})$。于是
\[
(a,b)=\bigcup_{k=1}^{\infty}(a,b-1/k]\in\mathscr{B}(\mathbb{R}).
\]
由差集封闭性，
\[
\{b\}=(a,b]\setminus(a,b)\in\mathscr{B}(\mathbb{R}).
\]
同理 $\{a\}\in\mathscr{B}(\mathbb{R})$，再由有限并封闭性可得
\[
[a,b]=(a,b]\cup\{a\},
\qquad
[a,b)=(a,b)\cup\{a\}
\]
均属于 $\mathscr{B}(\mathbb{R})$。
\end{proof}

\subsection*{概率的公理化定义和性质}

概率的公理化定义继承了频率的非负性、规范性和可加性。

\begin{mydefinition}{概率与概率空间}{probability-space}
设 $\Omega$ 是非空集合，$\mathscr{F}$ 是 $\Omega$ 上给定的一个 $\sigma$-域。若映射
\[
P:\mathscr{F}\longrightarrow[0,+\infty)
\]
满足：
\begin{enumerate}
\item \textbf{非负性：}对任意 $A\in\mathscr{F}$，有 $P(A)\geq0$；
\item \textbf{规范性：}$P(\Omega)=1$；
\item \textbf{可列可加性：}若 $A_n\in\mathscr{F}$（$n\geq1$）两两互斥，则
\[
P\left(\bigcup_{n=1}^{\infty}A_n\right)
=\sum_{n=1}^{\infty}P(A_n),
\]
\end{enumerate}
则称 $P$ 为 $\mathscr{F}$ 上的\textbf{概率}，三元组 $(\Omega,\mathscr{F},P)$ 称为\textbf{概率空间}。
\end{mydefinition}
可以看到的是，$P$实际上就是一个测度，概率空间$(\Omega,\mathscr{F},P)$是一个测度空间.
\subsection*{概率的性质}
\begin{proposition}{空集的概率}{probability-empty-set}
\[
P(\varnothing)=0.
\]
\end{proposition}

\begin{proof}
令 $A_1=\Omega$，$A_n=\varnothing$（$n\geq2$）。则 $A_1,A_2,\ldots$ 两两互斥，且其并为 $\Omega$。由规范性和可列可加性，
\[
1=P(\Omega)=\sum_{n=1}^{\infty}P(A_n)
=P(\Omega)+\sum_{n=2}^{\infty}P(\varnothing).
\]
若 $P(\varnothing)>0$，右端将大于 $1$，矛盾，故 $P(\varnothing)=0$。
\end{proof}

\textbf{注：}当 $A=\varnothing$ 时，必有 $P(A)=0$；反过来，$P(A)=0$ 并不能推出 $A=\varnothing$。例如在几何概型 $\Omega=[0,1]$ 中，单点集 $\{a\}$（$a\in[0,1]$）满足
\[
P(\{a\})=\dfrac{m(\{a\})}{m(\Omega)}=0,
\]
但 $\{a\}$ 不是空集。满足 $P(A)=0$ 的事件称为\textbf{零概率集}。

\begin{proposition}{有限可加性}{finite-additivity}
若 $A_1,\ldots,A_n\in\mathscr{F}$ 两两互斥，则
\[
P\left(\bigcup_{i=1}^{n}A_i\right)
=\sum_{i=1}^{n}P(A_i).
\]
\end{proposition}

\begin{proof}
令 $A_{n+k}=\varnothing$（$k\geq1$），将有限个事件补成可列个两两互斥事件。由可列可加性及 $P(\varnothing)=0$，即得结论。
\end{proof}

\begin{proposition}{单调性}{probability-monotonicity}
若 $A,B\in\mathscr{F}$ 且 $B\subseteq A$，则
\[
P(B)\leq P(A).
\]
\end{proposition}

\begin{proof}
因为 $A=B\cup(A\setminus B)$，且两部分互斥，故
\[
P(A)=P(B)+P(A\setminus B)\geq P(B).
\]
\end{proof}

\begin{mycorollary}{概率的上界}{probability-upper-bound}
对任意 $A\in\mathscr{F}$，有
\[
0\leq P(A)\leq1.
\]
\end{mycorollary}

\begin{proof}
左侧由非负性成立；又 $A\subseteq\Omega$，故由单调性和规范性，$P(A)\leq P(\Omega)=1$。
\end{proof}

\begin{proposition}{补事件的概率}{probability-complement}
对任意 $A\in\mathscr{F}$，有
\[
P(A^c)=1-P(A).
\]
\end{proposition}

\begin{proof}
由于 $A\cap A^c=\varnothing$ 且 $A\cup A^c=\Omega$，由有限可加性得
\[
P(A)+P(A^c)=P(\Omega)=1.
\]
\end{proof}

\begin{proposition}{一般加法公式}{probability-addition-formula}
对任意事件 $A,B\in\mathscr{F}$，有
\[
P(A\cup B)=P(A)+P(B)-P(A\cap B).
\]
\end{proposition}

\begin{proof}
因为 $A\cup B=A\cup(B\setminus A)$，且 $A$ 与 $B\setminus A$ 互斥，故
\[
P(A\cup B)=P(A)+P(B\setminus A).
\]
又 $B=(B\setminus A)\cup(A\cap B)$ 为互斥并，故
\[
P(B\setminus A)=P(B)-P(A\cap B).
\]
代入即得结论。
\end{proof}

\textbf{注：}三个事件 $A,B,C$ 的一般加法公式为
\[
\begin{aligned}
P(A\cup B\cup C)={}&P(A)+P(B)+P(C)\\
&-P(A\cap B)-P(A\cap C)-P(B\cap C)\\
&+P(A\cap B\cap C).
\end{aligned}
\]

\begin{mytheorem}{容斥恒等式}{inclusion-exclusion}
对任意 $n$ 个事件 $A_1,A_2,\ldots,A_n$，有
\[
P\left(\bigcup_{i=1}^{n}A_i\right)
=\sum_{r=1}^{n}(-1)^{r+1}
\sum_{1\leq i_1<\cdots<i_r\leq n}
P\left(A_{i_1}\cap\cdots\cap A_{i_r}\right).
\]
\end{mytheorem}

\begin{proof}
对 $n$ 归纳。$n=1$ 时结论显然成立。设结论对 $n-1$ 个事件成立，令 $U_{n-1}=\bigcup_{i=1}^{n-1}A_i$。由一般加法公式，
\[
P(U_{n-1}\cup A_n)=P(U_{n-1})+P(A_n)-P(U_{n-1}\cap A_n).
\]
又
\[
U_{n-1}\cap A_n=\bigcup_{i=1}^{n-1}(A_i\cap A_n).
\]
分别对 $U_{n-1}$ 与右侧这个 $n-1$ 个事件的并使用归纳假设，合并不含 $A_n$ 和含 $A_n$ 的各项，便得到 $n$ 个事件的公式。
\end{proof}

\begin{mycorollary}{有限次次可加性}{finite-subadditivity}
对任意 $n$ 个事件 $A_1,\ldots,A_n$，有
\[
P\left(\bigcup_{i=1}^{n}A_i\right)
\leq\sum_{i=1}^{n}P(A_i).
\]
\end{mycorollary}

\begin{proof}
对 $n$ 归纳。$n=1$ 时结论显然成立。若结论对 $n-1$ 个事件成立，则由一般加法公式和 $P\bigl((\bigcup_{i=1}^{n-1}A_i)\cap A_n\bigr)\geq0$，有
\[
\begin{aligned}
P\left(\bigcup_{i=1}^{n}A_i\right)
&\leq P\left(\bigcup_{i=1}^{n-1}A_i\right)+P(A_n)\\
&\leq\sum_{i=1}^{n}P(A_i).
\end{aligned}
\]
\end{proof}

\begin{example}[Kounias 不等式]
对任意 $n$ 个事件 $A_1,\ldots,A_n$，有
\[
P\left(\bigcup_{r=1}^{n}A_r\right)
\leq\min_{1\leq k\leq n}
\left\{
\sum_{r=1}^{n}P(A_r)-\sum_{\substack{1\leq r\leq n\\ r\ne k}}P(A_r\cap A_k)
\right\}.
\]
\end{example}

\begin{proof}
固定 $k$。由有限次次可加性，
\[
\begin{aligned}
\sum_{r=1}^{n}P(A_r)-\sum_{r\ne k}P(A_r\cap A_k)
&=P(A_k)+\sum_{r\ne k}P(A_r\cap A_k^c)\\
&\geq P\left(A_k\cup\bigcup_{r\ne k}(A_r\cap A_k^c)\right)\\
&=P\left(\bigcup_{r=1}^{n}A_r\right).
\end{aligned}
\]
对 $k=1,\ldots,n$ 取最小值即得结论。
\end{proof}

\begin{mydefinition}{事件序列的极限}{limit-of-event-sequences}
若事件列 $\{A_n\}_{n\geq1}$ 满足
\[
A_1\subseteq A_2\subseteq\cdots\subseteq A_n\subseteq\cdots,
\]
则称其为\textbf{递增序列}（或不降序列），并定义
\[
\lim_{n\to\infty}A_n:=\bigcup_{n=1}^{\infty}A_n.
\]
若
\[
A_1\supseteq A_2\supseteq\cdots\supseteq A_n\supseteq\cdots,
\]
则称其为\textbf{递减序列}（或不增序列），并定义
\[
\lim_{n\to\infty}A_n:=\bigcap_{n=1}^{\infty}A_n.
\]
\end{mydefinition}

\begin{mytheorem}{概率的连续性}{continuity-of-probability}
\begin{enumerate}
\item 若 $\{A_n\}_{n\geq1}$ 为递增序列，则
\[
P\left(\lim_{n\to\infty}A_n\right)=\lim_{n\to\infty}P(A_n).
\]
称为概率的\textbf{下连续性}；
\item 若 $\{A_n\}_{n\geq1}$ 为递减序列，则
\[
P\left(\lim_{n\to\infty}A_n\right)=\lim_{n\to\infty}P(A_n).
\]
称为概率的\textbf{上连续性}。
\end{enumerate}
\end{mytheorem}

\begin{proof}
对递增序列，令 $F_1=A_1$，$F_n=A_n\setminus A_{n-1}$（$n\geq2$）。则 $F_1,F_2,\ldots$ 两两互斥，且
\[
\bigcup_{i=1}^{n}F_i=A_n,
\qquad
\bigcup_{i=1}^{\infty}F_i=\bigcup_{i=1}^{\infty}A_i.
\]
因此
\[
P\left(\lim_{n\to\infty}A_n\right)
=\sum_{i=1}^{\infty}P(F_i)
=\lim_{n\to\infty}\sum_{i=1}^{n}P(F_i)
=\lim_{n\to\infty}P(A_n).
\]
对递减序列，$\{A_n^c\}_{n\geq1}$ 为递增序列。由补事件公式和已证的下连续性，
\[
\begin{aligned}
P\left(\bigcap_{n=1}^{\infty}A_n\right)
&=1-P\left(\bigcup_{n=1}^{\infty}A_n^c\right)\\
&=1-\lim_{n\to\infty}P(A_n^c)
=\lim_{n\to\infty}P(A_n).
\end{aligned}
\]
\end{proof}

\begin{proposition}{可列次可加性（布尔不等式）}{countable-subadditivity}
对任意可列个事件 $A_1,A_2,\ldots$，有
\[
P\left(\bigcup_{i=1}^{\infty}A_i\right)
\leq\sum_{i=1}^{\infty}P(A_i).
\]
\end{proposition}

\begin{proof}
令 $B_n=\bigcup_{i=1}^{n}A_i$。则 $\{B_n\}$ 为递增序列，且
\[
\lim_{n\to\infty}B_n=\bigcup_{i=1}^{\infty}A_i.
\]
由下连续性和有限次次可加性，
\[
\begin{aligned}
P\left(\bigcup_{i=1}^{\infty}A_i\right)
&=\lim_{n\to\infty}P(B_n)\\
&\leq\lim_{n\to\infty}\sum_{i=1}^{n}P(A_i)
=\sum_{i=1}^{\infty}P(A_i).
\end{aligned}
\]
\end{proof}

\begin{example}[可列个零概率事件的并]
若事件列 $\{A_i\}_{i\geq1}$ 满足 $P(A_i)=0$（$i\geq1$），则
\[
P\left(\bigcup_{i=1}^{\infty}A_i\right)=0.
\]
\end{example}

\begin{proof}
由布尔不等式，
\[
0\leq P\left(\bigcup_{i=1}^{\infty}A_i\right)
\leq\sum_{i=1}^{\infty}P(A_i)=0.
\]
\end{proof}

\textbf{注：}若事件列 $\{A_i\}_{i\geq1}$ 满足 $P(A_i)=1$（$i\geq1$），则 $P(A_i^c)=0$。因此
\[
P\left(\bigcap_{i=1}^{\infty}A_i\right)
=1-P\left(\bigcup_{i=1}^{\infty}A_i^c\right)=1.
\]

\begin{mydefinition}{事件列的上极限与下极限}{limsup-liminf-events}
设 $A_n\in\mathscr{F}$（$n\geq1$），不要求事件列 $\{A_n\}$ 单调。定义
\[
\limsup_{n\to\infty}A_n
:=\bigcap_{n=1}^{\infty}\bigcup_{k=n}^{\infty}A_k
\]
为事件列的\textbf{上极限}，定义
\[
\liminf_{n\to\infty}A_n
:=\bigcup_{n=1}^{\infty}\bigcap_{k=n}^{\infty}A_k
\]
为事件列的\textbf{下极限}。
\end{mydefinition}

\textbf{注：}对任意样本点 $\omega$，有
\[
\begin{aligned}
\limsup_{n\to\infty}A_n\text{发生}&\iff\omega\in\limsup_{n\to\infty}A_n\\&\iff \omega\in \bigcap_{n=1}^{\infty}\bigcup_{k=n}^{\infty}A_k,\\
&\iff \forall n\geqslant 1,\omega \in \bigcup_{k=n}^{\infty}A_k\\&\iff \forall n\geqslant 1,\exists k_n\geqslant n ,\text{使得}\omega\in A_{k_n}\\
&\iff \text{存在无穷子序列}A_{k_n}, k_n\geqslant n,\text{使得},\omega \in A_{k_n}\\
& \iff \text{存在无穷多个事件}A_{k_n},k_n\geqslant 1,\text{发生}
\end{aligned}
\]
\[
\begin{aligned}
	\liminf_{n\to\infty}A_n\text{发生}&\iff\omega\in\liminf_{n\to\infty}A_n\\&\iff \omega\in \bigcup_{n=1}^{\infty}\bigcap_{k=n}^{\infty}A_k,\\
	&\iff \exists n\geqslant 1,\omega \in \bigcap_{k=n}^{\infty}A_k\\&\iff \exists n\geqslant 1,\text{使得}\forall k\geqslant n ,\text{有}\omega\in A_{k}\\
	&\iff\exists n\geqslant 1,\text{使得}\forall k\geqslant n ,\text{有}A_{k}\text{发生}\\
	& \iff \text{最多只有有限个时间}A_k\text{不发生}
\end{aligned}
\]
因此，上极限表示“无穷多个 $A_n$ 发生”，下极限表示“只有有限个 $A_n$ 不发生”，并且恒有
\[
\liminf_{n\to\infty}A_n\subseteq\limsup_{n\to\infty}A_n.
\]
对于一般的事件列的极限，我们可以利用上下极限给出定义。
\begin{mydefinition}
	当一个事件列的上极限与下极限相等时，也即
	\[\liminf_{n\to\infty}A_n =\limsup_{n\to\infty}A_n\]
	时，我们定义这个事件列的极限为
	\[\lim\limits_{n\to\infty}A_n :=\liminf_{n\to\infty}A_n =\limsup_{n\to\infty}A_n\]
\end{mydefinition}
\subsection*{一些概率空间的例子和计算}

\begin{example}[伯努利概率空间]
设 $\Omega$ 是非空集合，$\varnothing\subsetneq A\subsetneq\Omega$，并令
\[
\mathscr{F}=\{\varnothing,\Omega,A,A^c\}.
\]
给定 $0<p<1$，定义
\[
P(\varnothing)=0,\qquad P(\Omega)=1,
\qquad P(A)=p,\qquad P(A^c)=1-p.
\]
则 $(\Omega,\mathscr{F},P)$ 构成一个概率空间，称为\textbf{伯努利概率空间}（或伯努利概率模型），它用于讨论单个事件是否发生。
\end{example}

\begin{example}[古典概率空间]
设 $\Omega$ 只含有限个样本点，并取
\[
\mathscr{F}=2^\Omega=\{A:A\subseteq\Omega\}.
\]
若 $n(A)$ 表示集合 $A$ 中元素的个数，定义
\[
P(A)=\frac{n(A)}{n(\Omega)},\qquad A\in\mathscr{F},
\]
则 $(\Omega,\mathscr{F},P)$ 构成古典概率空间。这里每个单点事件的概率均为 $1/n(\Omega)$。
\end{example}

\textbf{注：}计算古典概型时，可以不逐项写出样本空间和 $\sigma$-域，但必须先确认样本空间中的各单点事件等可能。

\begin{example}[最大号码问题]
某城市有编号为 $1,2,\ldots,N$ 的 $N$ 部卡车。独立记录遇到的 $n$ 辆卡车的编号，允许重复。求记录到的最大编号为 $k$（$1\leq k\leq N$）的概率。
\end{example}

\begin{proof}
令 $B_k$ 表示最大编号恰为 $k$，$A_k$ 表示最大编号不超过 $k$。则
\[
B_1=A_1,\qquad B_k=A_k\setminus A_{k-1}\quad(2\leq k\leq N),
\]
且
\[
P(A_k)=\left(\frac{k}{N}\right)^n.
\]
因此
\[
P(B_1)=\frac{1}{N^n},
\]
而当 $2\leq k\leq N$ 时，
\[
P(B_k)=P(A_k)-P(A_{k-1})
=\frac{k^n-(k-1)^n}{N^n}.
\]

\end{proof}

\begin{example}[配对问题]
在有 $n$ 人参加的晚会上，每人带来一顶各不相同的帽子。离场时，每人从随机混放的帽子中抽取一顶。求至少有一人抽到自己帽子的概率。
\end{example}

\begin{proof}
给人和帽子分别编号 $1,2,\ldots,n$，令 $A_i$ 表示第 $i$ 人抽到自己的帽子。对任意
$1\leq i_1<i_2<\cdots<i_r\leq n$，有
\[
P(A_{i_1}\cap\cdots\cap A_{i_r})=\frac{(n-r)!}{n!}.
\]
同阶交事件共有 $\binom{n}{r}$ 个，故
\[
\sum_{1\leq i_1<\cdots<i_r\leq n}
P(A_{i_1}\cap\cdots\cap A_{i_r})
=\binom{n}{r}\frac{(n-r)!}{n!}=\frac{1}{r!}.
\]
由容斥恒等式，所求概率为
\[
P\left(\bigcup_{i=1}^{n}A_i\right)
=\sum_{r=1}^{n}\frac{(-1)^{r+1}}{r!}
=1-\sum_{r=0}^{n}\frac{(-1)^r}{r!}.
\]
因此
\[
\lim_{n\to\infty}P\left(\bigcup_{i=1}^{n}A_i\right)
=1-e^{-1}\approx0.6321,
\]
而无人抽到自己帽子的概率趋于 $e^{-1}\approx0.3679$。
\end{proof}

\begin{example}[圆桌夫妻问题]
十对夫妻随机地沿圆桌坐下，求没有妻子坐在丈夫身边的概率。
\end{example}

\begin{proof}
令 $A_i$ 表示第 $i$ 对夫妻相邻（$i=1,2,\ldots,10$）。总坐法数为
$(20-1)!=19!$。指定 $r$ 对夫妻相邻时，将每对夫妻看作一个整体，并计入夫妻内部的两种次序，故
\[
P(A_{i_1}\cap\cdots\cap A_{i_r})
=\frac{2^r(20-r-1)!}{19!}
=\frac{2^r(19-r)!}{19!}.
\]
由容斥恒等式，至少有一对夫妻相邻的概率为
\[
P\left(\bigcup_{i=1}^{10}A_i\right)
=\sum_{r=1}^{10}(-1)^{r+1}\binom{10}{r}
\frac{2^r(19-r)!}{19!}
\approx0.6605.
\]
所以，没有夫妻相邻的概率约为
\[
1-0.6605=0.3395.
\]
\end{proof}

\begin{example}[反复抽取扑克牌中的大王]
从一副 $54$ 张的扑克牌中反复抽牌，每次抽取一张，记录后放回。求：
\begin{enumerate}
\item 一直抽不到大王的概率；
\item 无穷多次抽到大王的概率。
\end{enumerate}
\end{example}

\begin{proof}
设 $E_n$ 表示第 $n$ 次抽到大王。

对第一个问题，令
\[
C_n=\bigcap_{k=1}^{n}E_k^c.
\]
则 $C_n$ 表示前 $n$ 次均未抽到大王，且 $\{C_n\}$ 为递减序列。令
\[
C=\bigcap_{n=1}^{\infty}C_n,
\]
则 $C$ 表示始终抽不到大王。由独立性和概率的上连续性，
\[
P(C)=\lim_{n\to\infty}P(C_n)
=\lim_{n\to\infty}\left(\frac{53}{54}\right)^n=0.
\]
因此，最终至少抽到一次大王的概率为 $1$。

对第二个问题，令 $B$ 表示无穷多次抽到大王，则
\[
B=\limsup_{n\to\infty}E_n
=\bigcap_{n=1}^{\infty}\bigcup_{k=n}^{\infty}E_k.
\]
再令
\[
B_n=\bigcup_{k=n}^{\infty}E_k.
\]
事件 $B_n$ 表示“从第 $n$ 次抽牌起至少抽到一次大王”。由第一个问题同样的计算，
$P(B_n^c)=0$，故 $P(B_n)=1$。由于 $\{B_n\}$ 为递减序列，
\[
P(B)=P\left(\bigcap_{n=1}^{\infty}B_n\right)
=\lim_{n\to\infty}P(B_n)=1.
\]
\end{proof}

\textbf{注：}本例中“最终抽到大王”和“无穷多次抽到大王”的概率均为 $1$，但“概率为 $1$”并不等同于逻辑上的必然发生。把概率小于 $0.05$ 的事件称为“小概率事件”也是常见约定，而不是普遍的数学定义。

\subsection*{测度及其基本性质}

概率可以看作总质量为 $1$ 的测度。先将概率公理中的“规范性”放宽，即得到更一般的测度概念。

\begin{mydefinition}{测度、有限测度与 $\sigma$ 有限测度}{measure-space}
设 $\Omega$ 是非空集合，$\mathscr{F}$ 是 $\Omega$ 上给定的一个 $\sigma$-域。若映射
\[
\mu:\mathscr{F}\longrightarrow[0,+\infty]
\]
满足：
\begin{enumerate}
\item $\mu(\varnothing)=0$；
\item 对任意两两互斥的事件列 $A_1,A_2,\ldots\in\mathscr{F}$，有
\[
\mu\left(\bigcup_{n=1}^{\infty}A_n\right)
=\sum_{n=1}^{\infty}\mu(A_n),
\]
\end{enumerate}
则称 $\mu$ 为 $\mathscr{F}$ 上的\textbf{测度}，三元组 $(\Omega,\mathscr{F},\mu)$ 称为\textbf{测度空间}。

\end{mydefinition}
若存在两两互斥的事件列 $A_1,A_2,\ldots\in\mathscr{F}$，使得
\[
\bigcup_{n=1}^{\infty}A_n=\Omega,
\qquad
\mu(A_n)<+\infty\quad(n\geq1),
\]
则称 $\mu$ 为\textbf{$\sigma$ 有限测度}。若 $\mu(\Omega)<+\infty$，则称 $\mu$ 为\textbf{有限测度}；若 $\mu(\Omega)=1$，则称 $\mu$ 为\textbf{概率测度}，也简称为\textbf{概率}。
\textbf{注：}当 $\mu(\Omega)=1$ 时，上述定义恰好就是前面概率 $P$ 的公理化定义。因此，每个概率空间都是测度空间。

\begin{proposition}{测度的有限可加性}{measure-finite-additivity}
若 $A_1,\ldots,A_n\in\mathscr{F}$ 两两互斥，则
\[
\mu\left(\bigcup_{i=1}^{n}A_i\right)=\sum_{i=1}^{n}\mu(A_i).
\]
\end{proposition}

\begin{proof}
令 $A_{n+k}=\varnothing$（$k\geq1$），再对事件列 $A_1,A_2,\ldots$ 使用测度的可列可加性即可。
\end{proof}

\begin{proposition}{测度的单调性}{measure-monotonicity}
若 $A,B\in\mathscr{F}$ 且 $B\subseteq A$，则
\[
\mu(B)\leq\mu(A).
\]
\end{proposition}

\begin{proof}
由 $A=B\cup(A\setminus B)$，且右端两事件互斥，得
\[
\mu(A)=\mu(B)+\mu(A\setminus B)\geq\mu(B).
\]
\end{proof}

\begin{proposition}{测度的可减性}{measure-subtractivity}
若 $A,B\in\mathscr{F}$，$A\subseteq B$，且 $\mu(A)<+\infty$，则
\[
\mu(B\setminus A)=\mu(B)-\mu(A).
\]
\end{proposition}

\begin{proof}
由 $B=A\cup(B\setminus A)$，且右端两事件互斥，故
\[
\mu(B)=\mu(A)+\mu(B\setminus A).
\]
由于 $\mu(A)<+\infty$，可移项得到结论。
\end{proof}

\begin{proposition}{测度的容斥恒等式}{measure-inclusion-exclusion}
若 $A_1,A_2,\ldots,A_n\in\mathscr{F}$，且对每个 $1\leq k\leq n$ 都有 $\mu(A_k)<+\infty$，则
\[
\mu\left(\bigcup_{k=1}^{n}A_k\right)
=\sum_{r=1}^{n}(-1)^{r+1}
\sum_{1\leq i_1<i_2<\cdots<i_r\leq n}
\mu(A_{i_1}\cap A_{i_2}\cap\cdots\cap A_{i_r}).
\]
\end{proposition}

\begin{proof}
将概率容斥公式的证明中 $P$ 全部替换为 $\mu$ 即可；有限性条件保证各项的减法有意义。
\end{proof}

\begin{mytheorem}{测度的可列次可加性}{measure-countable-subadditivity}
对任意事件列 $A_1,A_2,\ldots\in\mathscr{F}$，有
\[
\mu\left(\bigcup_{n=1}^{\infty}A_n\right)
\leq\sum_{n=1}^{\infty}\mu(A_n).
\]
\end{mytheorem}

\begin{proof}
令 $B_1=A_1$，并对 $n\geq2$ 令
\[
B_n=A_n\setminus\bigcup_{k=1}^{n-1}A_k.
\]
则 $B_1,B_2,\ldots$ 两两互斥，$B_n\subseteq A_n$，且
\[
\bigcup_{n=1}^{\infty}A_n=\bigcup_{n=1}^{\infty}B_n.
\]
故由可列可加性和单调性，
\[
\mu\left(\bigcup_{n=1}^{\infty}A_n\right)
=\sum_{n=1}^{\infty}\mu(B_n)
\leq\sum_{n=1}^{\infty}\mu(A_n).
\]
\end{proof}

\begin{proposition}{测度的下连续性}{measure-lower-continuity}
若 $\{A_n\}_{n\geq1}\subseteq\mathscr{F}$ 为递增序列，则
\[
\mu\left(\lim_{n\to\infty}A_n\right)
=\lim_{n\to\infty}\mu(A_n).
\]
\end{proposition}

\begin{proof}
令 $F_1=A_1$，$F_n=A_n\setminus A_{n-1}$（$n\geq2$）。仿照概率下连续性的证明，对两两互斥的 $F_n$ 应用测度的可列可加性即可。
\end{proof}

\begin{proposition}{测度的上连续性}{measure-upper-continuity}
若 $\{A_n\}_{n\geq1}\subseteq\mathscr{F}$ 为递减序列，且 $\mu(A_1)<+\infty$，则
\[
\mu\left(\lim_{n\to\infty}A_n\right)
=\lim_{n\to\infty}\mu(A_n).
\]
\end{proposition}

\begin{proof}
令 $C_n=A_1\setminus A_n$。则 $\{C_n\}$ 为递增序列，且
\[
\bigcup_{n=1}^{\infty}C_n
=A_1\setminus\bigcap_{n=1}^{\infty}A_n.
\]
由下连续性和可减性，
\[
\begin{aligned}
\mu(A_1)-\mu\left(\bigcap_{n=1}^{\infty}A_n\right)
&=\mu\left(\bigcup_{n=1}^{\infty}C_n\right)\\
&=\lim_{n\to\infty}\mu(C_n)\\
&=\lim_{n\to\infty}\bigl[\mu(A_1)-\mu(A_n)\bigr].
\end{aligned}
\]
由于 $\mu(A_1)<+\infty$，整理即得结论。
\end{proof}

\begin{example}[Lebesgue 测度]
在 $\mathscr{B}(\mathbb{R}^d)$ 上存在唯一的测度 $\lambda$，使得对任意半开长方体
\[
(a_1,b_1]\times(a_2,b_2]\times\cdots\times(a_d,b_d]
\]
（其中 $a_i\leq b_i$，$1\leq i\leq d$），都有
\[
\lambda\left((a_1,b_1]\times\cdots\times(a_d,b_d]\right)
=(b_1-a_1)(b_2-a_2)\cdots(b_d-a_d).
\]
该测度称为 $\mathbb{R}^d$ 上的\textbf{Lebesgue 测度}，它给出了通常所说的长度、面积和体积。
\end{example}

\section{条件概率}

\subsection*{条件概率的定义}

在已知事件 $B$ 发生的条件下，原样本空间缩小为 $B$。因此，事件 $A$ 在这一条件下发生的可能性，应由 $A$ 与 $B$ 同时发生的概率相对于 $B$ 的概率来刻画。

\begin{example}[掷两枚骰子]
依次掷两枚均匀骰子。已知第一枚骰子的点数为 $4$，求两枚骰子点数之和为偶数的概率。

设 $B$ 表示“第一枚骰子的点数为 $4$”，$A$ 表示“两枚骰子点数之和为偶数”。在 $B$ 已经发生时，第二枚骰子共有 $6$ 种等可能结果；其中点数为 $2,4,6$ 时 $A$ 发生，故
\[
P(A\mid B)=\frac{3}{6}=\frac12.
\]
\end{example}

\begin{mydefinition}{条件概率}{conditional-probability}
设 $(\Omega,\mathscr{F},P)$ 为概率空间，$A,B\in\mathscr{F}$，且 $P(B)>0$。称
\[
P(A\mid B)=\frac{P(A\cap B)}{P(B)}
\]
为在事件 $B$ 发生的条件下事件 $A$ 的\textbf{条件概率}。
\end{mydefinition}

\begin{example}[寻找钥匙]
一串钥匙可能落在左口袋或右口袋。设 $L$ 表示“钥匙在左口袋”，$R$ 表示“钥匙在右口袋”，且
\[
P(L)=0.4,\qquad P(R)=0.4,\qquad R\subseteq L^{\mathrm c}.
\]
检查左口袋后发现没有钥匙，即已知 $L^{\mathrm c}$ 发生。于是
\[
P(R\mid L^{\mathrm c})
=\frac{P(R\cap L^{\mathrm c})}{P(L^{\mathrm c})}
=\frac{P(R)}{1-P(L)}
=\frac{0.4}{0.6}
=\frac23.
\]
\end{example}

\subsection*{乘法公式}

\begin{mytheorem}{条件概率的乘法公式}{conditional-multiplication-formula}
若 $A,B\in\mathscr{F}$ 且 $P(B)>0$，则
\[
P(A\cap B)=P(B)P(A\mid B).
\]

更一般地，若 $A_1,\ldots,A_n\in\mathscr{F}$ 且
\[
P(A_1\cap\cdots\cap A_{n-1})>0,
\]
则
\[
\begin{aligned}
P(A_1\cap\cdots\cap A_n)
={}&P(A_1)P(A_2\mid A_1)P(A_3\mid A_1\cap A_2)\cdots\\
&\cdot P(A_n\mid A_1\cap\cdots\cap A_{n-1}).
\end{aligned}
\]
\end{mytheorem}

\begin{proof}
第一式由条件概率的定义直接得到。对一般情形逐次使用
\[
P(A_1\cap\cdots\cap A_k)
=P(A_1\cap\cdots\cap A_{k-1})
P(A_k\mid A_1\cap\cdots\cap A_{k-1})
\]
即可。
\end{proof}

\textbf{注：}\quad 交事件 $A\cap B$ 的概率同时受到 $P(B)$ 与 $P(A\mid B)$ 的影响。若 $P(B)=0$，初等定义中的 $P(A\mid B)$ 没有意义；但因为 $P(A\cap B)=0$，在只书写乘法公式时常约定乘积 $P(B)P(A\mid B)$ 为 $0$。

\subsection*{全概率公式}

\begin{mytheorem}{全概率公式}{total-probability-formula}
设 $A_1,A_2,\ldots$ 构成样本空间 $\Omega$ 的一个可列分割，即
\[
A_i\cap A_j=\varnothing\quad(i\ne j),
\qquad
\bigcup_{i=1}^{\infty}A_i=\Omega,
\]
并且 $P(A_i)>0$。则对任意 $B\in\mathscr{F}$，
\[
P(B)=\sum_{i=1}^{\infty}P(A_i)P(B\mid A_i).
\]
当分割只有有限个事件时，相应地使用有限求和。
\end{mytheorem}

\begin{proof}
由于
\[
B=B\cap\Omega
=\bigcup_{i=1}^{\infty}(B\cap A_i),
\]
且各个 $B\cap A_i$ 两两互斥，所以由可列可加性与乘法公式，
\[
P(B)=\sum_{i=1}^{\infty}P(B\cap A_i)
=\sum_{i=1}^{\infty}P(A_i)P(B\mid A_i).
\]
\end{proof}

\textbf{注：}\quad 若分割中含有零概率事件，可以将它们从求和中删去；这些事件对 $P(B)$ 的贡献均为 $0$，而 $P(B\mid A_i)$ 的初等比值定义本身并不存在。

\begin{example}[同卵双胞胎的比例]
设 $A$ 表示“双胞胎为同卵”，$B$ 表示“双胞胎性别相同”。已知同卵双胞胎性别必然相同，异卵双胞胎两人的性别可视为独立且各以 $1/2$ 的概率为男或女，所以
\[
P(B\mid A)=1,
\qquad
P(B\mid A^{\mathrm c})=\frac12.
\]
若总体中双胞胎性别相同的比例为 $P(B)=0.64$，记 $P(A)=x$。由全概率公式，
\[
0.64=x+\frac12(1-x),
\]
解得
\[
x=0.28.
\]
因此同卵双胞胎约占 $28\%$。
\end{example}

\begin{example}[Pólya 罐模型]
罐中原有 $a$ 个红球与 $b$ 个蓝球，其中 $a,b>0$。每次随机取出一个球，放回后再加入 $c$ 个同色球，其中 $c\geq0$。证明：任意第 $k$ 次取到红球的概率都等于
\[
\frac{a}{a+b}.
\]
\end{example}

\begin{proof}
对 $k$ 作归纳。$k=1$ 时结论显然成立。假设结论对第 $k-1$ 次抽取成立。以第一次抽取的颜色作分割：若第一次取到红球，则此后罐中有 $a+c$ 个红球与 $b$ 个蓝球；若第一次取到蓝球，则有 $a$ 个红球与 $b+c$ 个蓝球。由归纳假设和全概率公式，
\[
\begin{aligned}
P(\text{第 $k$ 次为红球})
&=\frac{a}{a+b}\cdot\frac{a+c}{a+b+c}
+\frac{b}{a+b}\cdot\frac{a}{a+b+c}\\
&=\frac{a}{a+b}.
\end{aligned}
\]
故结论对所有 $k\geq1$ 成立。
\end{proof}

\begin{example}[配对问题的递推与恰有 $k$ 个配对]
将 $n$ 封信随机放入写有相应地址的 $n$ 个信封中。令 $E_n$ 表示“没有一封信放入正确的信封”，并记
\[
q_n=P(E_n),\qquad q_0=1,\quad q_1=0.
\]
则对 $n\geq2$，
\[
q_n=\frac1n q_{n-2}+\frac{n-1}{n}q_{n-1}.
\]
从而
\[
q_n-q_{n-1}=-\frac1n(q_{n-1}-q_{n-2})=\frac{(-1)^n}{n!},
\]
于是
\[
q_n=\sum_{r=0}^{n}\frac{(-1)^r}{r!}.
\]

进一步，恰有 $k$ 封信放入正确信封的概率为
\[
\binom{n}{k}\frac{(n-k)!q_{n-k}}{n!}
=\frac1{k!}\sum_{r=0}^{n-k}\frac{(-1)^r}{r!},
\qquad 0\leq k\leq n.
\]
\end{example}

\begin{proof}
记 $D_n=n!q_n$ 为 $n$ 个元素的错排数。先指定第一封信放入哪一个错误的信封，共有 $n-1$ 种选择。若该信封对应的信又恰好放入第一封信的信封，剩余部分有 $D_{n-2}$ 种错排；否则把这两封信的放置关系合并后，可与 $n-1$ 个元素的错排一一对应，因而有 $D_{n-1}$ 种。因此
\[
D_n=(n-1)(D_{n-1}+D_{n-2}).
\]
两边除以 $n!$，便得到
\[
q_n=\frac1n q_{n-2}+\frac{n-1}{n}q_{n-1}.
\]

由 $q_1-q_0=-1$，反复使用差分递推即可得到 $q_n$ 的交错和。若恰有 $k$ 个位置正确，先选出这 $k$ 个位置，再将其余 $n-k$ 个位置完全错排，便得到最后一个公式。
\end{proof}

\subsection*{贝叶斯公式}

\begin{mytheorem}{贝叶斯公式}{bayes-formula}
设 $A_1,A_2,\ldots$ 构成 $\Omega$ 的一个可列分割，且 $P(A_i)>0$。若 $B\in\mathscr{F}$ 且 $P(B)>0$，则
\[
P(A_i\mid B)
=\frac{P(A_i)P(B\mid A_i)}
{\displaystyle\sum_{k=1}^{\infty}P(A_k)P(B\mid A_k)}.
\]
有限分割的情形只需将分母改为有限求和。
\end{mytheorem}

\begin{proof}
由条件概率的定义、乘法公式与全概率公式，
\[
P(A_i\mid B)
=\frac{P(A_i\cap B)}{P(B)}
=\frac{P(A_i)P(B\mid A_i)}
{\displaystyle\sum_{k=1}^{\infty}P(A_k)P(B\mid A_k)}.
\]
\end{proof}

\textbf{注：}\quad $P(A_i)$ 常称为事件 $A_i$ 的\textbf{先验概率}；获得信息 $B$ 后得到的 $P(A_i\mid B)$ 常称为\textbf{后验概率}。贝叶斯公式把结果出现后的信息用于修正原有判断。

\begin{example}[保险事故原因]
设 $A$ 表示“事故较容易发生”，$B$ 表示“发生事故”。已知
\[
P(A)=0.3,
\qquad
P(B\mid A)=0.4,
\qquad
P(B\mid A^{\mathrm c})=0.2.
\]
由全概率公式，
\[
P(B)=0.3\times0.4+0.7\times0.2=0.26.
\]
故
\[
P(A\mid B)
=\frac{0.3\times0.4}{0.26}
=\frac6{13}
\approx0.4615.
\]
\end{example}

\begin{example}[Monty Hall 问题]
三扇门后分别藏有一辆汽车和两只山羊。参赛者先选择第 $1$ 扇门；主持人知道汽车的位置，随后一定在第 $2,3$ 扇门中打开一扇藏有山羊的门。若两扇门都可打开，主持人以相同概率选择其中一扇。已知主持人打开了第 $3$ 扇门，问是否应改选第 $2$ 扇门。

令 $A_i$ 表示“汽车在第 $i$ 扇门后”，$B$ 表示“主持人打开第 $3$ 扇门”。则
\[
P(A_1)=P(A_2)=P(A_3)=\frac13,
\]
且
\[
P(B\mid A_1)=\frac12,
\qquad
P(B\mid A_2)=1,
\qquad
P(B\mid A_3)=0.
\]
因此
\[
P(B)=\frac13\cdot\frac12+\frac13\cdot1+\frac13\cdot0=\frac12.
\]
由贝叶斯公式，
\[
P(A_1\mid B)=\frac13,
\qquad
P(A_2\mid B)=\frac23.
\]
所以改选第 $2$ 扇门的获胜概率为 $2/3$，应当换门。
\end{example}

\subsection*{条件概率是一个概率}

\begin{mytheorem}{条件概率构成概率测度}{conditional-probability-measure}
固定 $B\in\mathscr{F}$ 且 $P(B)>0$，定义
\[
P_B(A)=P(A\mid B),\qquad A\in\mathscr{F}.
\]
则 $P_B$ 是可测空间 $(\Omega,\mathscr{F})$ 上的一个概率测度。
\end{mytheorem}

\begin{proof}
显然 $P_B(A)\geq0$，且
\[
P_B(\Omega)=\frac{P(\Omega\cap B)}{P(B)}=1.
\]
若 $A_1,A_2,\ldots$ 两两互斥，则 $A_i\cap B$ 也两两互斥，因而
\[
\begin{aligned}
P_B\left(\bigcup_{i=1}^{\infty}A_i\right)
&=\frac{P\left(\bigcup_{i=1}^{\infty}(A_i\cap B)\right)}{P(B)}\\
&=\sum_{i=1}^{\infty}\frac{P(A_i\cap B)}{P(B)}
=\sum_{i=1}^{\infty}P_B(A_i).
\end{aligned}
\]
故 $P_B$ 满足概率公理。
\end{proof}

\textbf{注：}\quad 因为 $P_B$ 本身是概率，所以前面关于概率的性质、乘法公式、全概率公式和贝叶斯公式都可以在 $P_B$ 下使用。

\begin{mydefinition}{条件概率空间中的再次条件化}{iterated-conditioning}
固定 $B\in\mathscr{F}$ 且 $P(B)>0$。若 $C\in\mathscr{F}$ 且 $P_B(C)>0$，定义
\[
P_B(A\mid C)=\frac{P_B(A\cap C)}{P_B(C)}.
\]
\end{mydefinition}

\begin{proposition}{再次条件化等价于对交事件条件化}{nested-conditioning}
固定 $B\in\mathscr{F}$ 且 $P(B)>0$。对任意 $C\in\mathscr{F}$，有
\[
P_B(C)>0\quad\Longleftrightarrow\quad P(B\cap C)>0.
\]
在上述条件成立时，对任意 $A\in\mathscr{F}$，
\[
P_B(A\mid C)=P(A\mid B\cap C).
\]
\end{proposition}

\begin{proof}
由定义，
\[
P_B(C)=\frac{P(B\cap C)}{P(B)},
\]
而 $P(B)>0$，故第一个结论成立。进一步，
\[
\begin{aligned}
P_B(A\mid C)
&=\frac{P_B(A\cap C)}{P_B(C)}
=\frac{P(A\cap C\cap B)/P(B)}{P(C\cap B)/P(B)}\\
&=\frac{P(A\cap B\cap C)}{P(B\cap C)}
=P(A\mid B\cap C).
\end{aligned}
\]
\end{proof}

\textbf{注：}\quad 在所有条件概率都有定义时，把通常的乘法公式、全概率公式和贝叶斯公式应用于概率 $P_B$，可分别得到
\[
P(A\cap C\mid B)
=P(A\mid B)P(C\mid A\cap B),
\]
\[
P(C\mid B)
=\sum_i P(A_i\mid B)P(C\mid A_i\cap B),
\]
以及
\[
P(A_j\mid B\cap C)
=\frac{P(A_j\mid B)P(C\mid A_j\cap B)}
{\displaystyle\sum_i P(A_i\mid B)P(C\mid A_i\cap B)}.
\]
这里第二、三式中的 $\{A_i\}$ 是 $\Omega$ 的一个有限或可列分割；求和只取使相关条件概率有定义的项，且贝叶斯公式的分母须为正。

\begin{example}[血液化验与重复检测]
某项血液化验对患者有 $95\%$ 的概率检出阳性，但对健康人也有 $1\%$ 的概率出现伪阳性。已知该疾病患者占总人口的 $0.5\%$。若某人的一次化验结果为阳性，求此人确实患病的条件概率。

设 $A$ 表示“此人患病”，$B$ 表示“第一次化验呈阳性”。则
\[
P(A)=0.005,
\qquad
P(B\mid A)=0.95,
\qquad
P(B\mid A^{\mathrm c})=0.01.
\]
由贝叶斯公式，
\[
\begin{aligned}
P(A\mid B)
&=\frac{P(A)P(B\mid A)}
{P(A)P(B\mid A)+P(A^{\mathrm c})P(B\mid A^{\mathrm c})}\\
&=\frac{0.005\times0.95}
{0.005\times0.95+(1-0.005)\times0.01}\\
&\approx0.323.
\end{aligned}
\]
因此，即使检验对患者的阳性检出率很高，一次阳性结果所对应的患病概率也只有约 $32.3\%$。原因在于患病率很低，健康人数量很大，伪阳性在全部阳性结果中占据了不可忽略的比例。

若此人再次接受同样的化验，记第二次阳性为 $B_2$。假设两次化验在给定患病状态后条件独立，即
\[
P(B_2\mid A\cap B)=P(B_2\mid A)=0.95,
\qquad
P(B_2\mid A^{\mathrm c}\cap B)=P(B_2\mid A^{\mathrm c})=0.01,
\]
则连续两次呈阳性后，
\[
\begin{aligned}
P(A\mid B\cap B_2)
&=\frac{P(A)P(B\mid A)P(B_2\mid A)}
{P(A)P(B\mid A)P(B_2\mid A)
+P(A^{\mathrm c})P(B\mid A^{\mathrm c})P(B_2\mid A^{\mathrm c})}\\
&=\frac{0.005\times0.95^2}
{0.005\times0.95^2+0.995\times0.01^2}\\
&\approx0.978.
\end{aligned}
\]
连续两次阳性使后验患病概率提高到约 $97.8\%$。
\end{example}

\begin{example}[飞机失事搜索]
一架飞机失踪，并且等可能地落在第 $1,2,3$ 号区域。设 $B_i$ 表示“飞机落在第 $i$ 个区域”，$A_i$ 表示“搜索第 $i$ 个区域后没有找到飞机”，并设
\[
P(A_i\mid B_i)=\beta_i,
\qquad 0<\beta_i<1,
\qquad i=1,2,3.
\]
若飞机不在被搜索的区域，则在那里搜索必然找不到飞机。

搜索第 $1$ 号区域而没有找到飞机后，由贝叶斯公式，
\[
P(B_1\mid A_1)=\frac{\beta_1}{\beta_1+2},
\qquad
P(B_2\mid A_1)=P(B_3\mid A_1)=\frac{1}{\beta_1+2}.
\]

随后又搜索第 $2$ 号区域，仍未找到飞机，则
\[
\begin{aligned}
P(B_1\mid A_1\cap A_2)&=\frac{\beta_1}{\beta_1+\beta_2+1},\\
P(B_2\mid A_1\cap A_2)&=\frac{\beta_2}{\beta_1+\beta_2+1},\\
P(B_3\mid A_1\cap A_2)&=\frac{1}{\beta_1+\beta_2+1}.
\end{aligned}
\]

若三个区域依次搜索后都没有找到飞机，则
\[
P(B_i\mid A_1\cap A_2\cap A_3)
=\frac{\beta_i}{\beta_1+\beta_2+\beta_3},
\qquad i=1,2,3.
\]
可见每次搜索失败后，都应使用新的后验概率作为下一次更新的先验概率。
\end{example}

\begin{example}[赌徒破产模型]
设 $a$ 为正整数，质点从整数点 $i\in\{0,1,\ldots,a\}$ 出发。只要质点位于 $1,2,\ldots,a-1$ 中，它就以概率 $p$ 向右移动一步，以概率 $q=1-p$ 向左移动一步，其中 $0<p<1$；一旦到达 $0$ 或 $a$，便停留在那里。求质点最终到达 $0$ 与到达 $a$ 的概率。

设 $E_i$ 表示“质点从 $i$ 出发”，$F$ 表示“最终到达 $0$”，并记
\[
Q_i=P(F\mid E_i).
\]
按第一次移动的方向使用全概率公式，得到
\[
Q_i=pQ_{i+1}+qQ_{i-1},
\qquad 1\leq i\leq a-1,
\]
边界条件为
\[
Q_0=1,
\qquad
Q_a=0.
\]
递推式可改写为
\[
Q_{i+1}-Q_i=\frac{q}{p}(Q_i-Q_{i-1}).
\]

当 $p=q=1/2$ 时，各阶差分相等，结合边界条件可得
\[
Q_i=\frac{a-i}{a}.
\]
当 $p\ne q$ 时，对等比差分求和并代入边界条件，得到
\[
Q_i
=\frac{\left(\dfrac qp\right)^i-\left(\dfrac qp\right)^a}
{1-\left(\dfrac qp\right)^a}.
\]
因此
\[
P(F\mid E_i)=
\begin{cases}
\dfrac{\left(\dfrac qp\right)^i-\left(\dfrac qp\right)^a}
{1-\left(\dfrac qp\right)^a},&p\ne q,\\[1.2em]
\dfrac{a-i}{a},&p=q=\dfrac12.
\end{cases}
\]

同理，若 $M$ 表示“最终到达 $a$”，则
\[
P(M\mid E_i)=
\begin{cases}
\dfrac{1-\left(\dfrac qp\right)^i}
{1-\left(\dfrac qp\right)^a},&p\ne q,\\[1.2em]
\dfrac{i}{a},&p=q=\dfrac12.
\end{cases}
\]
两式相加得到
\[
P(F\cup M\mid E_i)=P(F\mid E_i)+P(M\mid E_i)=1,
\]
即质点以概率 $1$ 被某一个端点吸收。

这个模型也可解释为赌徒破产问题：甲、乙的初始赌本分别为 $i$ 元和 $a-i$ 元，每局胜者从败者处获得 $1$ 元，甲每局获胜的概率为 $p$。固定甲的初始赌本 $i$，令乙的初始赌本趋于无穷，即令 $a\to\infty$，则
\[
\lim_{a\to\infty}P(F\mid E_i)=
\begin{cases}
\left(\dfrac qp\right)^i,&p>\dfrac12,\\[0.8em]
1,&p\leq\dfrac12.
\end{cases}
\]
因此，当 $p>1/2$ 时，甲有正概率永不破产；当 $p\leq1/2$ 时，面对拥有无限赌资的对手，甲以概率 $1$ 破产。
\end{example}

\textbf{注：}\quad 上述推导假设 $0<p<1$。退化情形 $p=0$ 或 $p=1$ 可直接判断：除非一开始已经位于相反的吸收端点，质点将确定地分别到达 $0$ 或 $a$。“以概率 $1$ 破产”表示几乎必然破产，并不把概率为 $1$ 与逻辑上的绝对必然混为一谈。

\section{事件的独立性和试验的独立性}

\subsection*{独立事件}

通常，事件 $B$ 的发生会改变事件 $A$ 发生的可能性，即 $P(A\mid B)\ne P(A)$。在某些情形下，$B$ 的发生不影响 $A$ 的概率，此时由乘法公式可得到
\[
P(A\cap B)=P(A)P(B).
\]
后一等式不要求 $P(B)>0$，因此适合作为独立性的定义。

\begin{mydefinition}{两个事件的独立性}{independence-two-events}
设 $(\Omega,\mathscr{F},P)$ 为概率空间，$A,B\in\mathscr{F}$。若
\[
P(A\cap B)=P(A)P(B),
\]
则称事件 $A$ 与 $B$\textbf{相互独立}。
\end{mydefinition}

\textbf{注：}\quad 当 $P(B)>0$ 时，
\[
P(A\mid B)=P(A)
\quad\Longleftrightarrow\quad
P(A\cap B)=P(A)P(B).
\]
独立性的定义本身仍可用于 $P(B)=0$ 的情形。

\begin{proposition}{独立性对补事件保持}{independence-complements}
若事件 $A$ 与 $B$ 相互独立，则下面三对事件也分别相互独立：
\[
A\text{ 与 }B^{\mathrm c},
\qquad
A^{\mathrm c}\text{ 与 }B,
\qquad
A^{\mathrm c}\text{ 与 }B^{\mathrm c}.
\]
\end{proposition}

\begin{proof}
由 $A$ 与 $B$ 独立，
\[
\begin{aligned}
P(A\cap B^{\mathrm c})
&=P(A)-P(A\cap B)\\
&=P(A)-P(A)P(B)\\
&=P(A)P(B^{\mathrm c}),
\end{aligned}
\]
故 $A$ 与 $B^{\mathrm c}$ 独立。交换 $A,B$ 的位置可得 $A^{\mathrm c}$ 与 $B$ 独立。最后，
\[
\begin{aligned}
P(A^{\mathrm c}\cap B^{\mathrm c})
&=1-P(A)-P(B)+P(A\cap B)\\
&=(1-P(A))(1-P(B))\\
&=P(A^{\mathrm c})P(B^{\mathrm c}),
\end{aligned}
\]
所以 $A^{\mathrm c}$ 与 $B^{\mathrm c}$ 独立。
\end{proof}

\textbf{注：}\quad 相互独立与互斥是两个不同的概念。互斥要求 $P(A\cap B)=0$，独立要求 $P(A\cap B)=P(A)P(B)$。若两个事件既互斥又独立，则必有 $P(A)=0$ 或 $P(B)=0$；因此两个概率均为正的互斥事件不可能独立。

\begin{example}[独立与互斥的区别]
向区间 $[0,1]$ 内均匀随机投入一个粒子。

令
\[
A_1=\left[0,\frac12\right],
\qquad
B_1=\left[\frac14,\frac34\right].
\]
则
\[
A_1\cap B_1=\left[\frac14,\frac12\right],
\]
并且
\[
P(A_1)=\frac12,
\qquad
P(B_1)=\frac12,
\qquad
P(A_1\cap B_1)=\frac14=P(A_1)P(B_1).
\]
所以 $A_1$ 与 $B_1$ 不互斥，但相互独立。

再令
\[
A_2=\left[\frac15,\frac13\right],
\qquad
B_2=\left[\frac14,1\right].
\]
则
\[
A_2\cap B_2=\left[\frac14,\frac13\right],
\]
且
\[
P(A_2)=\frac{2}{15},
\qquad
P(B_2)=\frac34,
\qquad
P(A_2\cap B_2)=\frac1{12}\ne\frac1{10}=P(A_2)P(B_2).
\]
所以 $A_2$ 与 $B_2$ 既不互斥，也不相互独立。
\end{example}
\begin{mydefinition}{两两独立与相互独立}{pairwise-mutual-independence}
	设 $A_1,A_2,\ldots,A_n$ 为 $n\geq2$ 个事件。
	
	若对任意 $i\ne j$ 都有
	\[
	P(A_i\cap A_j)=P(A_i)P(A_j),
	\]
	则称 $A_1,A_2,\ldots,A_n$\textbf{两两独立}。
	
	若对任意 $k\in\{2,3,\ldots,n\}$ 以及任意不同的指标
	\[
	1\leq i_1<i_2<\cdots<i_k\leq n,
	\]
	都有
	\[
	P(A_{i_1}\cap A_{i_2}\cap\cdots\cap A_{i_k})
	=\prod_{r=1}^{k}P(A_{i_r}),
	\]
	则称 $A_1,A_2,\ldots,A_n$\textbf{相互独立}。
\end{mydefinition}

\textbf{注：}\quad 相互独立一定推出两两独立；当 $n\geq3$ 时，两两独立一般不能推出相互独立。
\begin{example}[两两独立不推出相互独立]
掷两枚均匀骰子。令 $E$ 表示“两枚骰子的点数之和为 $7$”，$F$ 表示“第一枚骰子的点数为 $4$”，$G$ 表示“第二枚骰子的点数为 $3$”。则
\[
P(E)=P(F)=P(G)=\frac16,
\]
并且
\[
P(E\cap F)=P(E\cap G)=P(F\cap G)=P(E\cap F\cap G)=\frac1{36}.
\]
因此
\[
P(E\cap F)=P(E)P(F),
\qquad
P(E\cap G)=P(E)P(G),
\qquad
P(F\cap G)=P(F)P(G),
\]
所以 $E,F,G$ 两两独立。但是
\[
P(E\cap F\cap G)=\frac1{36}
\ne\frac1{216}=P(E)P(F)P(G),
\]
故它们并不相互独立。
\end{example}

\begin{example}[三色球模型]
袋中有 $4$ 个球：一个只涂白色，一个只涂蓝色，一个只涂红色，另一个同时涂有白、蓝、红三色。随机取出一个球。设 $A,B,C$ 分别表示“取出的球涂有白色、蓝色、红色”。则
\[
P(A)=P(B)=P(C)=\frac12,
\]
且
\[
P(A\cap B)=P(A\cap C)=P(B\cap C)=P(A\cap B\cap C)=\frac14.
\]
因此
\[
P(A\cap B)=P(A)P(B),
\quad
P(A\cap C)=P(A)P(C),
\quad
P(B\cap C)=P(B)P(C),
\]
所以 $A,B,C$ 两两独立；但
\[
P(A\cap B\cap C)=\frac14
\ne\frac18=P(A)P(B)P(C),
\]
故它们不相互独立。
\end{example}

\textbf{注：}\quad 按定义逐项验证 $n$ 个事件相互独立，需要检查
\[
\binom n2+\binom n3+\cdots+\binom nn=2^n-n-1
\]
个等式。在有明确应用背景时，通常先依据试验机制判断独立性，再用独立性完成计算。

\begin{mydefinition}{可列个事件的相互独立}{countable-event-independence}
设 $A_1,A_2,\ldots$ 为可列个事件。若其中任意有限多个不同事件都相互独立，则称事件列 $\{A_n\}_{n\geq1}$\textbf{相互独立}。
\end{mydefinition}

\subsection*{独立试验}

\begin{mydefinition}{独立试验、伯努利试验与独立重复试验}{independent-trials}
设一个试验由一系列子试验组成。
\begin{enumerate}
    \item 若分别由不同子试验决定的任意有限个事件相互独立，则称这些子试验\textbf{相互独立}。
    \item 若一个试验由 $n$ 个相互独立且同分布的子试验组成，则称其为\textbf{$n$ 重独立试验}。若每个子试验只有“成功”和“失败”两种结果，且成功概率为 $p$，$0<p<1$，则称其为\textbf{$n$ 重伯努利试验}。
    \item 若一个试验由可列个相互独立且同分布的子试验组成，则称其为\textbf{独立重复试验}。
\end{enumerate}
\end{mydefinition}
更加严格的定义看下一节.

\begin{example}[伯努利试验中的常用计算]
进行独立重复伯努利试验，每次成功的概率为 $p$，失败的概率为 $q=1-p$。
\begin{enumerate}
    \item 前 $n$ 次中恰好成功 $k$ 次的概率为
    \[
    \binom nk p^kq^{n-k}.
    \]
    \item 第一次成功发生在第 $k$ 次的概率为
    \[
    q^{k-1}p.
    \]
    \item 第 $r$ 次成功发生在第 $k$ 次的概率为
    \[
    \binom{k-1}{r-1}p^rq^{k-r},
    \qquad 1\leq r\leq k.
    \]
    \item 令 $F_n$ 表示前 $n$ 次全部失败，则
    \[
    F_n\downarrow\bigcap_{n=1}^{\infty}F_n
    =\{\text{永远不成功}\}.
    \]
    由概率的上连续性，
    \[
    P(\text{永远不成功})
    =\lim_{n\to\infty}P(F_n)
    =\lim_{n\to\infty}q^n=0.
    \]
    因而最终至少成功一次的概率为 $1$。
    \item “第 $n$ 次成功出现在第 $m$ 次失败之前”的概率可写成三种等价形式：
    \[
    \sum_{k=0}^{m-1}\binom{k+n-1}{n-1}p^nq^k,
    \]
    \[
    \sum_{k=n}^{\infty}\binom{k+m-1}{k}p^kq^m,
    \]
    或
    \[
    \sum_{k=n}^{m+n-1}\binom{m+n-1}{k}p^kq^{m+n-1-k}.
    \]
    三式分别按第 $n$ 次成功前的失败次数、第 $m$ 次失败前的成功次数，以及前 $m+n-1$ 次中的成功总数分类。
\end{enumerate}
\end{example}

\textbf{注：}\quad “小概率事件终会发生”只有在固定成功概率 $p>0$、各次试验相互独立并且试验无限重复等条件下才成立，不能脱离条件作为一般原则使用。概率为 $1$ 也不等同于逻辑上的绝对必然。

\begin{example}[分赌注问题]
甲、乙约定先胜 $t$ 局者获得全部赌注。比赛在甲已胜 $r$ 局、乙已胜 $s$ 局时中止，其中 $r<t$、$s<t$。若以后每局甲获胜的概率为 $p$，且各局相互独立，则
\[
\begin{aligned}
P(\text{甲最终先胜 }t\text{ 局})
&=P(t-r\text{ 次成功发生在 }t-s\text{ 次失败之前})\\
&=\sum_{k=0}^{t-s-1}
\binom{k+t-r-1}{t-r-1}p^{t-r}(1-p)^k.
\end{aligned}
\]
合理的分配方式是按甲、乙最终先胜 $t$ 局的概率分配赌注。
\end{example}

\begin{example}[事件 $A$ 先于事件 $B$ 发生]
进行独立重复试验。每次试验中，互斥事件 $A,B$ 发生的概率分别为 $p_A,p_B$，并设 $p_A+p_B>0$。于是
\[
\begin{aligned}
P(A\text{ 先于 }B\text{ 发生})
&=\sum_{n=1}^{\infty}(1-p_A-p_B)^{n-1}p_A\\
&=\frac{p_A}{p_A+p_B}.
\end{aligned}
\]
也可以令所求概率为 $x$。按第一次试验分类，
\[
x=p_A+(1-p_A-p_B)x,
\]
解得同一结果。
\end{example}

\begin{example}[由独立性得到欧拉乘积公式]
设 $s>1$，在正整数集合 $\mathbb Z_{>0}$ 上定义
\[
P(A)=\frac{\displaystyle\sum_{k\in A}k^{-s}}
{\displaystyle\sum_{k=1}^{\infty}k^{-s}}.
\]
将所有素数依次记为 $p_1,p_2,\ldots$，令
\[
A_i=\{mp_i:m\in\mathbb Z_{>0}\}.
\]
则 $P(A_i)=p_i^{-s}$。任取不同的 $i_1,\ldots,i_r$，有
\[
A_{i_1}\cap\cdots\cap A_{i_r}
=\{mp_{i_1}\cdots p_{i_r}:m\in\mathbb Z_{>0}\},
\]
所以
\[
P(A_{i_1}\cap\cdots\cap A_{i_r})
=\frac1{p_{i_1}^s\cdots p_{i_r}^s}
=\prod_{j=1}^{r}P(A_{i_j}).
\]
因此 $A_1,A_2,\ldots$ 相互独立。

又因为每个大于 $1$ 的正整数至少有一个素因子，
\[
\{1\}=\bigcap_{i=1}^{\infty}A_i^{\mathrm c}.
\]
由概率的上连续性和独立性，
\[
P(\{1\})
=\prod_{i=1}^{\infty}\left(1-\frac1{p_i^s}\right).
\]
另一方面，
\[
P(\{1\})=\left(\sum_{k=1}^{\infty}\frac1{k^s}\right)^{-1}.
\]
故得到欧拉乘积公式
\[
\sum_{k=1}^{\infty}\frac1{k^s}
=\prod_{i=1}^{\infty}\left(1-\frac1{p_i^s}\right)^{-1}.
\]
\end{example}
\newpage
\section{关于试验独立性的补充}

\subsection*{乘积概率空间}

设子试验 $1$ 和子试验 $2$ 的概率空间分别为
\[
(\Omega_1,\mathscr F_1,P_1),
\qquad
(\Omega_2,\mathscr F_2,P_2).
\]
两个子试验组成的大试验以
\[
\Omega=\Omega_1\times\Omega_2
\]
为样本空间。子试验 $1$ 中的事件 $A\in\mathscr F_1$ 在大样本空间中表示为 $A\times\Omega_2$；子试验 $2$ 中的事件 $B\in\mathscr F_2$ 表示为 $\Omega_1\times B$。两者同时发生对应
\[
(A\times\Omega_2)\cap(\Omega_1\times B)=A\times B.
\]

\begin{example}[掷骰子与独立取球]
子试验 $1$ 为掷一枚均匀骰子，$\Omega_1=\{1,2,3,4,5,6\}$；子试验 $2$ 为从红、白、蓝三个球中等概率取一个，$\Omega_2=\{\text{红},\text{白},\text{蓝}\}$。若两个子试验相互独立，则每个联合结果 $(i,c)\in\Omega_1\times\Omega_2$ 的概率为
\[
P(\{(i,c)\})=\frac16\cdot\frac13=\frac1{18}.
\]
更一般地，对任意 $A\subseteq\Omega_1$、$B\subseteq\Omega_2$，有
\[
P(A\times B)=P_1(A)P_2(B).
\]
\end{example}

\begin{mydefinition}{两个独立子试验的乘积空间表述}{two-independent-experiments-product}
在 $\Omega_1\times\Omega_2$ 上取乘积 $\sigma$ 域
\[
\mathscr F_1\otimes\mathscr F_2
=\sigma\bigl(\{A\times B:A\in\mathscr F_1,\ B\in\mathscr F_2\}\bigr).
\]
若其上的概率 $P$ 对任意 $A\in\mathscr F_1$、$B\in\mathscr F_2$ 都满足
\[
P(A\times B)=P_1(A)P_2(B),
\]
则称这两个子试验相互独立，并将 $P$ 记为 $P_1\otimes P_2$。
\end{mydefinition}

若 $P_2(B)>0$，则
\[
\begin{aligned}
P(A\times\Omega_2\mid\Omega_1\times B)
&=\frac{P(A\times B)}{P(\Omega_1\times B)}\\
&=\frac{P_1(A)P_2(B)}{P_1(\Omega_1)P_2(B)}\\
&=P_1(A).
\end{aligned}
\]
这严格表达了“子试验 $2$ 中事件 $B$ 的发生不影响子试验 $1$ 中事件 $A$ 的概率”。

\begin{mydefinition}{有限个独立子试验}{finite-independent-experiments-product}
设第 $i$ 个子试验的概率空间为 $(\Omega_i,\mathscr F_i,P_i)$，$i=1,\ldots,n$。在
\[
\Omega=\prod_{i=1}^{n}\Omega_i,
\qquad
\mathscr F=\bigotimes_{i=1}^{n}\mathscr F_i
\]
上，若概率 $P$ 对任意 $A_i\in\mathscr F_i$ 都满足
\[
P(A_1\times\cdots\times A_n)
=\prod_{i=1}^{n}P_i(A_i),
\]
则称这 $n$ 个子试验相互独立，并记
\[
P=P_1\otimes\cdots\otimes P_n.
\]
若各个概率空间完全相同，则它们组成 $n$ 重独立试验。
\end{mydefinition}

\subsection*{不独立子试验与转移概率}

\begin{example}[掷骰子后按点数取球]
先掷一枚均匀骰子。若点数为偶数，则从红、白两球中等概率取一个；若点数为奇数，则取蓝球。用
\[
p(\omega_1,\{\omega_2\})
=P(\text{第二步取出 }\omega_2\mid\text{第一步点数为 }\omega_1)
\]
表示转移概率。若 $\omega_1=2,4,6$，则
\[
p(\omega_1,\{\text{红}\})=p(\omega_1,\{\text{白}\})=\frac12,
\qquad
p(\omega_1,\{\text{蓝}\})=0;
\]
若 $\omega_1=1,3,5$，则
\[
p(\omega_1,\{\text{红}\})=p(\omega_1,\{\text{白}\})=0,
\qquad
p(\omega_1,\{\text{蓝}\})=1.
\]

对事件 $A\subseteq\Omega_1$ 和颜色集合 $B\subseteq\Omega_2$，联合概率可写为
\[
P(A\times B)
=\sum_{\omega_1\in A}P_1(\{\omega_1\})
\sum_{\omega_2\in B}p(\omega_1,\{\omega_2\}).
\]
例如，
\[
P(\text{取出红球})=\frac14,
\qquad
P(\text{骰子点数为 }2)=\frac16,
\]
而
\[
P(\text{点数为 }2\text{ 且取出红球})=\frac1{12}
\ne\frac16\cdot\frac14.
\]
所以两个子试验不独立。
\end{example}

